% ============================================================
% chapter_99_appendices.tex - Appendices A through H
%
% Purpose: archival back-matter material kept separate from the body
%          chapters: data sources, code+data availability, full
%          configuration template, per-station verification table,
%          reproducibility checklist, glossary, and the spatial atlas.
% Sections: A Data Catalogue, B Code and Data Availability,
%           C Configuration Template, D Per-Station Verification,
%           E Reproducibility Checklist, F Glossary, G Spatial Atlas,
%           H Supplementary Temporal-Alignment Diagnostics.
% Numbering: \thefigure / \thetable / \theequation prefix re-set to
%            the appendix letter so figures land as G.1, G.2, ...
%            instead of clashing with the Chapter 4 numbering.
% Included by main_thesis.tex (after Chapter 6).
% License: CC BY 4.0 (see ../LICENSE).
% ============================================================

% ----- APPENDICES divider: title at top of page -----
\clearpage
\thispagestyle{fancy}
\begin{center}
  {\fontsize{14}{17}\selectfont\bfseries APPENDICES}
\end{center}
\addcontentsline{toc}{chapter}{APPENDICES}
\vspace{\baselineskip}

\appendix
% In the appendix the chapter counter restarts; the body-chapter numbering
% scheme would produce table/figure numbers like "4.1" (chapter 4 ==
% Appendix D), colliding with Chapter 4. Switch the prefix to the appendix
% letter so the first table in Appendix D becomes "D.1" and the list of
% tables/figures shows letters rather than digits.
\renewcommand{\thetable}{\Alph{chapter}.\arabic{table}}
\renewcommand{\thefigure}{\Alph{chapter}.\arabic{figure}}
\renewcommand{\theequation}{\Alph{chapter}.\arabic{equation}}
% Make the appendices flow continuously: remove the forced page break that
% \chapter normally issues, so each appendix follows the previous one on the
% same page (a new page is taken only when the content runs over). This patch
% is applied here, after the main chapters, so only appendices A-F are
% affected; the biography uses its own explicit \clearpage.
\makeatletter
\patchcmd{\chapter}{\if@openright\cleardoublepage\else\clearpage\fi}{}{}{}
\makeatother
% Appendix chapter headings: left-aligned, "A. DATA CATALOGUE" (heading 1).
\titleformat{\chapter}[hang]
  {\normalfont\fontsize{14}{17}\selectfont\bfseries}
  {\thechapter.}{0.6em}{\MakeUppercase}
\titlespacing*{\chapter}{0pt}{1.5\baselineskip}{6pt}

% ============================================================
\chapter{Data Catalogue}\label{app:data}
\addcontentsline{apx}{appx}{\protect\numberline{Appendix~\thechapter}Data Catalogue}
% ============================================================

\section{Datasets}

The study uses three precipitation datasets and one validation
network. Acronyms used across the appendices follow the same usage
as the main chapters and are summarised inline at first use here.
The Integrated Multi-satellitE Retrievals for Global Precipitation
Measurement (GPM) Late Run (IMERG-L, V07, $0.1^{\circ}$,
approximately 14-hour latency) is the satellite product corrected
\citep{ref:huffman2020}. The IMERG Final Run (IMERG-F, V07,
$0.1^{\circ}$) is retained as an independent benchmark but is not
used in the correction. The Climate Prediction Center (CPC)
Unified Gauge-Based Analysis of Global Daily Precipitation
(CPC-UNI, $0.5^{\circ}$) is the correction reference
\citep{ref:chen2008}. Satellite data were obtained from the
National Aeronautics and Space Administration (NASA) Goddard Earth
Sciences Data and Information Services Center (GES DISC); CPC-UNI
from the National Oceanic and Atmospheric Administration (NOAA)
Physical Sciences Laboratory. All fields span 2001 to 2025 and are
processed per dekad (days 1 to 10, 11 to 20, 21 to month end).

\section{BMKG validation network: 180 to 172}

Independent validation uses daily station observations from
Indonesia's Meteorological, Climatological, and Geophysical Agency
(BMKG). The parent network has approximately 180 stations; two
filters reduce it to the 172 used in validation:

\begin{enumerate}
    \item \textbf{Spatial-mapping filter.} The station coordinates
    must fall within the land mask used by the satellite product.
    A small number of stations near the maritime border or with
    mis-recorded coordinates fail this test.
    \item \textbf{Minimum paired-day filter.} After sentinel
    removal and pairing against the corrected satellite series,
    the station must retain at least \texttt{MIN\_VALID\_DAYS}
    (approximately 30) paired daily observations per dekadal
    window.
\end{enumerate}

Together these drop 8 of 180 stations, leaving 172 for the
per-station Taylor diagrams, multi-threshold verification, and
station-level summaries.

\section{Sentinel handling}

The BMKG archive encodes missing data with the sentinel value
$8888$. Across the raw record this sentinel occurs $76{,}778$
times ($8.02\%$ of raw cells); leaving it unfiltered would inflate
heavy-threshold event counts catastrophically (by more than
$1000\%$ at the $100$~mm/day threshold, because every sentinel
exceeds every physical threshold). The loader removes the sentinel
through the centralised \texttt{MISSING\_SENTINEL} configuration
value before any analysis, leaving $880{,}702$ valid daily
observations with a maximum physical daily total of $1{,}965.5$~mm.
All threshold-based verification in this thesis uses the
sentinel-filtered observations. A different station archive with a
different sentinel convention (for example $9999$ or $-999$) would
require the sentinel value to be reset in the configuration.


% ============================================================
\chapter{Code and Data Availability}\label{app:code}
\addcontentsline{apx}{appx}{\protect\numberline{Appendix~\thechapter}Code and Data Availability}
% ============================================================

The framework is released as an open implementation. The source
code, the six numbered notebooks, and the Quarto documentation
site are hosted on GitHub at
\url{https://github.com/bennyistanto/hybrid-bias-correction}
under the Mozilla Public License 2.0. The notebooks open directly
in Google Colab from the repository, and a province-scale Bali
example dataset is bundled in the repository so the pipeline can be
run without any external download.

The full Indonesia operational inputs (IMERG, CPC, and BMKG for
2001 to 2025, approximately $1.7$~GB) and the complete outputs
(corrected NetCDFs, metrics, quality assessment, station
validation, and figures, approximately $40$~GB) are deposited on
Zenodo with a citable DOI:
\url{https://doi.org/10.5281/zenodo.20287847}. To reproduce the
Indonesia results, a reader clones the repository, downloads the
Zenodo bundle, points \texttt{config.yml} at the extracted
directories, and runs the notebooks. The companion manuscript is
under review at \emph{Remote Sensing} (MDPI)\MDPIstatus; its DOI
and citation will be added on acceptance.


% ============================================================
\chapter{Configuration Template}\label{app:config}
\addcontentsline{apx}{appx}{\protect\numberline{Appendix~\thechapter}Configuration Template}
% ============================================================

The pipeline is driven by a single \texttt{config.yml} file that
resolves environment placeholders (\texttt{\{main\_dir\}},
\texttt{\{input\_dir\}}, \texttt{\{output\_dir\}}) at load time.
The annotated excerpt below shows the parameters most relevant to
this thesis; the full file is in the repository.

{\footnotesize
\begin{verbatim}
# --- Area of interest (Bali example subdomain) ---
aoi:
  lat_range: [-8.85, -8.05]
  lon_range: [114.45, 115.75]
# --- Distribution fitting ---
statistics:
  gpd_threshold_percentile: 80   # GPD tail threshold (wet-day pctl)
# --- Deep learning blend ---
deep_learning:
  blend_alpha: 0.70              # LSEQM weight at full confidence
  use_confidence_mask: true
# --- Station-density confidence mask ---
station_density:
  saturation_count: 2            # gauges per CPC cell for full confidence
  smoothing_sigma: 1.0
# --- Quality assessment (CQI component weights) ---
quality_assessment:
  component_weights:
    basic_stats: 0.4
    distribution: 0.3
    temporal: 0.3
\end{verbatim}
}

\enlargethispage{\baselineskip}%   keep the short closing paragraph on the same page as the verbatim
The three sensitivity parameters identified in
Chapter~\ref{chap:methods} (\texttt{blend\_alpha},
\texttt{gpd\_threshold\_percentile}, \texttt{saturation\_count})
are all set here, which is what makes the sensitivity sweep of
Section~\ref{sec:results-sensitivity} a configuration change 
rather than a code change.


% ============================================================
\chapter{Selected Per-Station Verification}\label{app:stations}
\addcontentsline{apx}{appx}{\protect\numberline{Appendix~\thechapter}Selected Per-Station Verification}
% ============================================================

Table~\ref{tab:perstation-bali} reports the per-station Taylor
statistics for the four BMKG stations on Bali, computed against
the daily station records over 2001 to 2021. The table illustrates
the correlation pattern at the level of individual stations: the
Pearson correlation barely moves between the raw IMERG-L and the
corrected LSEQM+DL series, while the standard deviation of the
corrected series moves toward the observed standard deviation. The
single-station daily samples are smaller and noisier than the
pooled 172-station Indonesia result, but the same pattern is
visible. The full per-station tables for all 172 Indonesia
stations are in the Zenodo release.
Table~\ref{tab:perstation-bali} reports, per station, the correlation
($r$) against the daily station record for raw IMERG-L and corrected
LSEQM+DL, the observed and LSEQM+DL standard deviations ($\sigma$, in
mm/day), and the LSEQM+DL bias (mm/day).

\begin{table}[H]
\centering
\caption{Per-station Taylor statistics for the four Bali BMKG stations.\label{tab:perstation-bali}}
\begingroup\small
\begin{tabular}{@{}lccccc@{}}
\toprule
Station (WMO ID) & $r$ IMERG-L & $r$ LSEQM+DL & $\sigma_{\text{obs}}$ & $\sigma_{\text{LSEQM+DL}}$ & Bias \\
\midrule
Ngurah Rai (97230)   & 0.387 & 0.388 & 14.08 & 14.32 & $+0.13$ \\
Denpasar (97232)     & 0.387 & 0.390 & 17.17 & 15.37 & $-0.73$ \\
Kahang-Kahang (97234)& 0.286 & 0.268 & 14.42 & 12.58 & $-0.47$ \\
Jembrana (97236)     & 0.233 & 0.202 & 17.36 & 11.52 & $-2.54$ \\
\bottomrule
\end{tabular}
\endgroup
\end{table}

\medskip
The single-month noise in such per-station figures is large enough to be
worth quantifying, because the optimal-window analysis of
Section~\ref{sec:results-window} reports per-station correlations, and an
earlier single-month version of that analysis produced systematically
larger per-station values. Figure~\ref{fig:noise-selection} explains the
difference. A per-station correlation estimated from a single month rests
on roughly ten to twenty wet days, so it scatters widely around its true
value: at a fixed aggregation window the single-month estimate has a
median near the seven-year value (both near $0.56$) but an
interquartile range of $[0.35, 0.73]$. Taking the maximum over the $97$
candidate window offsets then adds a selection bias, shifting the
single-month median up to $0.66$ and inflating the high tail, so that
half the stations appear to exceed $r = 0.65$. Both effects vanish on the
full seven-year record, where each offset is estimated precisely: the
per-station peak correlation tightens to a median of $0.56$ with an
interquartile range of $[0.51, 0.62]$, and the selection bias disappears
because the peak and fixed-window values coincide. The per-station
figures reported in the main text use the seven-year estimates for this
reason.

\begin{figure}[H]
\centering
\includegraphics[width=\textwidth]{fig_thesis_99_noise_selection.png}
\caption[Sampling stability of the per-station window correlation]{Sampling
stability of the per-station window correlation: one month is inflated by
noise and offset selection, the seven-year record is precise (dashed lines
mark each distribution's median).\label{fig:noise-selection}}
\end{figure}


% ============================================================
\chapter{Reproducibility Checklist}\label{app:repro}
\addcontentsline{apx}{appx}{\protect\numberline{Appendix~\thechapter}Reproducibility Checklist}
% ============================================================

This checklist records what an independent reader needs in order
to reproduce the results, following the reproducibility design of
Section~\ref{sec:methods-reproducibility} and the empirical
evaluation of Section~\ref{sec:results-reproducibility}.

\medskip\noindent\textbf{Datasets and access.}
IMERG-L and IMERG-F V07 from NASA GES DISC; CPC-UNI from NOAA PSL;
BMKG station records as bundled (Bali example) or via the Zenodo
deposit (full Indonesia). DOI:
\url{https://doi.org/10.5281/zenodo.20287847}.

\medskip\noindent\textbf{Software environment (verified run).}
Python $3.12$ on the Google Colab default image; TensorFlow
$2.20.0$; NumPy $2.0.2$; netCDF4 $1.7.4$; xarray, SciPy, and
scikit-learn from the Colab image. No pinned requirements file
beyond a single supplementary install was needed because the
remaining dependencies are pre-installed on Colab.

\medskip\noindent\textbf{Hardware tested.}
Free-tier Google Colab, standard CPU runtime with no hardware
accelerator; the two heaviest notebooks used the no-cost high-RAM
runtime option.

\medskip\noindent\textbf{Expected runtime.}
For the Bali subdomain ($126$-cell grid, $80$ land pixels) across
all $36$ dekadal windows, the end-to-end notebook sequence
(nb02 to nb06) completes in $72.1$ minutes
(Table~\ref{tab:colab-runtime}). The full Indonesian domain scales
up with pixel count and is correspondingly longer.

\medskip\noindent\textbf{Determinism.}
The statistical stages (LS, EQM, GPD) are deterministic given the
inputs. The CNN training uses a fixed random seed set in the
configuration; minor run-to-run variation at the level of
floating-point tolerance can remain from non-deterministic GPU or
threaded CPU operations, but the Bali run reported here used the
CPU runtime. A formal bit-equivalence assertion between the Colab
and local outputs is identified in
Section~\ref{sec:results-reproducibility} as the next
reproducibility deliverable.

\medskip\noindent\textbf{Run order.}
Notebook 01 (data acquisition or example bundle), then 02 (LSEQM+DL
correction), 03 (metrics), 04 (quality assessment), 05 (station
validation), 06 (visualisation). Only the setup cell of each
notebook and the upstream corrected outputs are prerequisites.


% ============================================================
\chapter{Glossary}\label{app:glossary}
\addcontentsline{apx}{appx}{\protect\numberline{Appendix~\thechapter}Glossary}
% ============================================================

% Glossary as a real two-column layout. Tunable widths:
%   labelwidth = width of the acronym (left) column
%   leftmargin = where the definition (right) column starts
%   (text width - leftmargin) = effective definition column width
% At 14 cm text width, leftmargin = 3.5 cm leaves 10.5 cm for definitions.
\begin{description}[noitemsep, topsep=2pt, align=left, style=multiline,
    labelindent=0pt, leftmargin=3.5cm, labelwidth=3.3cm, labelsep=0.2cm]
    \item[BMKG] Meteorological, Climatological, and Geophysical Agency.
    \item[BSQS] Basic Statistical Quality sub-score.
    \item[CNN] Convolutional Neural Network.
    \item[CPC-UNI] CPC Unified Gauge-Based Analysis of Global Daily Precipitation.
    \item[CQI] Composite Quality Index.
    \item[CRMSE] Centred root mean square error.
    \item[CSI] Critical Success Index.
    \item[DQS] Distribution Quality sub-score.
    \item[EQM] Empirical Quantile Mapping.
    \item[ETS] Equitable Threat Score.
    \item[FAR] False Alarm Ratio.
    \item[GPD] Generalized Pareto Distribution.
    \item[GPM] Global Precipitation Measurement mission.
    \item[IMERG] Integrated Multi-satellitE Retrievals for GPM.
    \item[IMERG-L/F ] IMERG Late Run / Final Run.
    \item[KS] Kolmogorov-Smirnov (test).
    \item[LS] Linear Scaling.
    \item[LSEQM+DL] The four-stage hybrid pipeline (LS, EQM, GPD, CNN).
    \item[MLE] Maximum likelihood estimation.
    \item[NSE] Nash-Sutcliffe Efficiency.
    \item[POD] Probability of Detection.
    \item[RB] Relative Bias.
    \item[RMSE] Root Mean Square Error.
    \item[SDR] Standard Deviation Ratio.
    \item[TQS] Temporal Quality sub-score.
    \item[WMO] World Meteorological Organization.
\end{description}


% ============================================================
\chapter{Spatial Atlas}\label{app:atlas}
\addcontentsline{apx}{appx}{\protect\numberline{Appendix~\thechapter}Spatial Atlas}
% ============================================================

% --- Tunable per-figure widths (fraction of \textwidth, 0.50 to 1.00) ---
% Adjust any of these to make a single figure smaller (more compact, leaves
% more vertical room on the page) or larger (more readable detail).
\newcommand{\atlasGoneWidth}{1.0}   % G.1 difference maps
\newcommand{\atlasGtwoWidth}{1.0}   % G.2 three-day x three-product
\newcommand{\atlasGthreeWidth}{1.0} % G.3 pillar metrics
\newcommand{\atlasGfourWidth}{0.95}  % G.4 six-product comparison
\newcommand{\atlasGfiveWidth}{0.95}  % G.5 station overlay

Single-day precipitation maps and per-pixel climatological metric maps for
the full Indonesian domain. The figures complement the summary statistics
of Chapter~\ref{chap:results} by giving a spatial view of the correction
at a representative day ($1$~January~$2025$), across the dekad-$1$-of-January
climatology where the per-pixel metrics are computed, and overlaying the
BMKG validation stations on the same maps. The representative day was
chosen because it falls within the active northwest-monsoon convective
regime over much of the archipelago, where the systematic bias of IMERG-L
against the gauge reference is the most visible.

Figure~\ref{fig:atlas-diff} maps three differences: the raw IMERG-L
minus CPC-UNI (where the satellite departs from the gauge reference),
the corrected LSEQM+DL minus CPC-UNI (the residual after correction),
and the corrected LSEQM+DL minus the raw IMERG-L (the net effect of the
pipeline, largest where the correction does the most work).

\begin{figure}[H]
\centering
\includegraphics[width=\atlasGoneWidth\textwidth]{fig_thesis_G1_differences.png}
\caption[Difference maps locating the correction residual]{Difference maps for $1$~January~$2025$ (mm/day).\label{fig:atlas-diff}}
\end{figure}

Figure~\ref{fig:atlas-threedays} contrasts three days (the start,
middle, and end of dekad~$1$: $1$, $5$, and $10$~January) across IMERG-L
(raw), CPC-UNI (gauge reference), and LSEQM+DL (corrected); the
correction tracks the day-to-day variation of the gauge reference at all
three regimes without imposing a uniform overlay, preserving spatial
pattern while bringing the daily amount toward the gauge target.

\begin{figure}[H]
\centering
\includegraphics[width=\atlasGtwoWidth\textwidth]{fig_thesis_G2_three_days.png}
\caption[Three-day daily contrast across products]{Three-day daily contrast across products for dekad~$1$ of January~$2025$.\label{fig:atlas-threedays}}
\end{figure}

Figure~\ref{fig:atlas-pillars} maps the four verification pillars
(rows: relative bias, KS $p$-value, $Q_{99}$ ratio, and CSI) across the
LS, LSEQM, and LSEQM+DL stages (columns), computed over all
dekad-$1$-of-January samples in the $2001$ to $2025$ record; each row's
colour scale is set so that the pillar's target value (zero for bias,
unity otherwise) reads as neutral colour.

\begin{figure}[H]
\centering
\includegraphics[width=\atlasGthreeWidth\textwidth]{fig_thesis_G3_pillar_metrics.png}
\caption[Per-pillar climatological metric maps across correction stages]{Per-pixel climatological metric maps for the four verification pillars across the three correction stages.\label{fig:atlas-pillars}}
\end{figure}

Figure~\ref{fig:atlas-six} maps six daily products on a shared colour
bar: the raw IMERG-L Late Run, the IMERG-F Final Run (the gauge-anchored
independent benchmark), the Linear Scaling stage, the LSEQM intermediate
(LS plus EQM plus GPD), the LSEQM+DL corrected output, and the CPC-UNI
gauge reference; reading the panels in sequence traces the correction
from raw retrieval to gauge reference.

\begin{figure}[H]
\centering
\includegraphics[width=\atlasGfourWidth\textwidth]{fig_thesis_G4_six_products.png}
\caption[Six-product comparison for 1 January 2025]{Six-product daily precipitation map for $1$~January~$2025$.\label{fig:atlas-six}}
\end{figure}

Figure~\ref{fig:atlas-stations} overlays the BMKG validation stations
(triangles coloured by their own dekad-$1$-of-January climatological
mean) on the same climatology for four products: CPC-UNI gauge analysis,
raw IMERG-L, corrected LSEQM+DL, and IMERG-F Final Run; where a
triangle's interior colour matches the surrounding gridded colour, the
gridded product agrees with the station.

\begin{figure}[H]
\centering
\includegraphics[width=\atlasGfiveWidth\textwidth]{fig_thesis_G5_station_overlay.png}
\caption[Station-overlay maps for dekad-1-January climatology]{Dekad-$1$-of-January climatological mean precipitation (mm/day, $2001$ to $2025$ average) for four products with the BMKG validation stations overlaid.\label{fig:atlas-stations}}
\end{figure}


% ============================================================
\chapter{Supplementary Temporal-Alignment Diagnostics}\label{app:window}
\addcontentsline{apx}{appx}{\protect\numberline{Appendix~\thechapter}Supplementary Temporal-Alignment Diagnostics}
% ============================================================

% --- Tunable per-figure widths (fraction of \textwidth) ---
\newcommand{\gridTrmmWidth}{1.0}   % H.1 gridded-CPC, TRMM-input era (full text width)
\newcommand{\gridGpmWidth}{1.0}    % H.2 gridded-CPC, GPM era (full text width)

The whole-domain window-offset diagnostic of
Section~\ref{sec:discussion-subdaily} (Figure~\ref{fig:gridded-cpc})
is shown here separately for the two satellite eras, with the
per-timezone-band pooled correlation the main-text figure omits.

\begin{figure}[H]
\centering
\includegraphics[width=\gridTrmmWidth\textwidth]{fig_thesis_H1_gridded_trmm.png}
\caption[Gridded window offset against CPC-UNI, TRMM-input era]{Gridded
window-offset diagnostic against CPC-UNI over the TRMM-input era ($2001$ to
$2013$).\label{fig:grid-trmm}}
\end{figure}

Each figure maps the optimal aggregation-window offset $h^\star$
of the half-hourly IMERG-L source against the gridded CPC-UNI
analysis at every $0.5^\circ$ cell, at CPC-UNI's native UTC labels
(panel a, $h^\star\approx0$) and after harmonising CPC-UNI to the
local-observation day (panel b, $h^\star\approx-23$~h), with the
pooled $r(h)$ and its per-band decomposition below (panel c). The
native optimum sits at $h^\star\approx0$ across the
gauge-constrained western and central cells in both eras, the
whole-domain signature of CPC-UNI sharing the IMERG-L UTC calendar
date; the data-sparse eastern cells, where the CPC-UNI gauge
network thins, carry low and noisy $h^\star$ and are read together
with CPC gauge density.

\begin{figure}[H]
\centering
\includegraphics[width=\gridGpmWidth\textwidth]{fig_thesis_H2_gridded_gpm.png}
\caption[Gridded window offset against CPC-UNI, GPM era]{Gridded window-offset
diagnostic against CPC-UNI over the GPM era ($2015$ to
$2025$).\label{fig:grid-gpm}}
\end{figure}

The optimum is the same in the timing-imprecise TRMM-input era
(Figure~\ref{fig:grid-trmm}) as in the GPM era
(Figure~\ref{fig:grid-gpm}), confirming that the offset is a
property of the gauge-labelling convention rather than of the
satellite retrieval skill. The absolute correlation against the
smoothed $0.5^\circ$ analysis is muted in both eras and so does
not reproduce the point-gauge era recovery documented in
Section~\ref{sec:results-window}, while the spatial coverage that
the gridded reference adds, beyond the BMKG station locations,
shows the offset to be uniform across the gauge-constrained part
of the archipelago.
